\subsection{Exact heat trace and the threshold-correct fixed-energy count}
\label{subsec:s6cusp-native-heat-count}

Retain the absolute ground energy
\begin{equation}\label{eq:s6cusp-native-ground-energy}
 E_{R;0,0}:=\mathcal B_{R,+}+\mathcal B_{R,-}
 =\frac{2\pi}{\delta_R}
 \sqrt{(p_R-6T_R)^2+144m_R^2}.
\end{equation}
The adjective ``ground'' in this definition is justified by the quadratic-
form identity \eqref{eq:s6cusp-native-oscillator-factorization}: for every
\(s\) in the form domain,
\begin{align}
 \left\langle s,
 (H_{\eta,R}-E_{R;0,0})s\right\rangle
 &=2\mathcal B_{R,+}\lVert\mathfrak a_+s\rVert^2
   +2\mathcal B_{R,-}\lVert\mathfrak a_-s\rVert^2
 \label{eq:s6cusp-native-ground-lower-bound}\\
 &\ge0.
\end{align}
Both coefficients are strictly positive.  Indeed,
\(\delta_R>0\) and
\(\Xi_R^2-S_R^2=24\delta_R\), where
\begin{equation}\label{eq:s6cusp-Xi-S-definitions}
 \Xi_R:=\sqrt{(p_R-6T_R)^2+144m_R^2},
 \qquad S_R:=p_R+6T_R,
\end{equation}
imply \(\Xi_R>|S_R|\).  Equations
\eqref{eq:s6cusp-native-frequency-plus}--
\eqref{eq:s6cusp-native-frequency-minus} therefore give
\(\mathcal B_{R,+},\mathcal B_{R,-}>0\).  The six-dimensional simultaneous
kernel \eqref{eq:s6cusp-native-joint-ground} proves that the lower bound is
attained with complex multiplicity six.  Thus topology supplies six states
\emph{at} \(E_{R;0,0}\); it supplies no state below that number.

\begin{theorem}[Exact native heat trace]
\label{thm:s6cusp-native-heat-trace}
For every \(t>0\),
\begin{equation}\label{eq:s6cusp-native-heat-trace}
 \operatorname{Tr}(e^{-tH_{\eta,R}})
 =\frac{3}{
  2\sinh(t\mathcal B_{R,+})
   \sinh(t\mathcal B_{R,-})}.
\end{equation}
At every fixed \(t>0\), the source radial cusp has the expansion
\begin{equation}\label{eq:s6cusp-native-heat-trace-cusp}
 \operatorname{Tr}(e^{-tH_{\eta,R}})
 =\frac{\delta_R}{16\pi^2t^2}
 \left[1+O_t(R^{-2})\right].
\end{equation}
\end{theorem}

\begin{proof}
The complete spectrum and its six-dimensional joint multiplicities give
\begin{align*}
 \operatorname{Tr}(e^{-tH_{\eta,R}})
 &=6\sum_{n_+,n_-\ge0}
 e^{-t[(2n_++1)\mathcal B_{R,+}
       +(2n_-+1)\mathcal B_{R,-}]}\\
 &=6\prod_{\sigma\in\{+,-\}}
 \frac{e^{-t\mathcal B_{R,\sigma}}}
 {1-e^{-2t\mathcal B_{R,\sigma}}}\\
 &=\frac{3}{2\sinh(t\mathcal B_{R,+})
                  \sinh(t\mathcal B_{R,-})},
\end{align*}
which proves \eqref{eq:s6cusp-native-heat-trace}.  The exact product
\eqref{eq:s6cusp-native-frequency-product} and
\(\mathcal B_{R,+},\mathcal B_{R,-}=O(R^{-1})\) give
\begin{align*}
 \sinh(t\mathcal B_{R,+})
 \sinh(t\mathcal B_{R,-})
 &=t^2\mathcal B_{R,+}\mathcal B_{R,-}
 \left[1+O_t(\mathcal B_{R,+}^2+
                  \mathcal B_{R,-}^2)\right]\\
 &=\frac{24\pi^2t^2}{\delta_R}
 \left[1+O_t(R^{-2})\right].
\end{align*}
Substitution into the exact trace proves
\eqref{eq:s6cusp-native-heat-trace-cusp}.
\end{proof}


For \(E>0\), define the absolute spectral counting function
\begin{equation}\label{eq:s6cusp-native-count-definition}
 N_{\eta,R}(E):=
 \operatorname{Tr}\mathbf1_{(-\infty,E]}(H_{\eta,R}).
\end{equation}
The exact lattice of admissible indices is
\begin{equation}\label{eq:s6cusp-native-count-index-set}
 \mathscr I_R(E):=
 \left\{(n_+,n_-)\in\mathbb Z_{\ge0}^2:
  2\mathcal B_{R,+}n_++2\mathcal B_{R,-}n_-
  \le E-E_{R;0,0}\right\}.
\end{equation}
It is empty when \(E<E_{R;0,0}\).  Consequently
\begin{equation}\label{eq:s6cusp-native-count-cardinality}
 N_{\eta,R}(E)=6\,\#\mathscr I_R(E).
\end{equation}
Writing
\begin{equation}\label{eq:s6cusp-native-count-L}
 L_R(E):=E-E_{R;0,0}
\end{equation}
without deleting its sign, the exact floor formula is
\begin{equation}\label{eq:s6cusp-native-count-piecewise}
 N_{\eta,R}(E)=
 \begin{cases}
 0,&E<E_{R;0,0},\\[0.4em]
 6\displaystyle\sum_{n_+=0}^{
   \left\lfloor L_R(E)/(2\mathcal B_{R,+})\right\rfloor}
 \left[
  1+\left\lfloor
  \dfrac{L_R(E)-2\mathcal B_{R,+}n_+}
        {2\mathcal B_{R,-}}
  \right\rfloor
 \right],&E\ge E_{R;0,0}.
 \end{cases}
\end{equation}
In particular,
\begin{equation}\label{eq:s6cusp-native-count-threshold-values}
 \operatorname{Tr}\mathbf1_{(-\infty,E_{R;0,0})}(H_{\eta,R})=0,
 \qquad
 N_{\eta,R}(E_{R;0,0})=6.
\end{equation}
Replacing \(L_R(E)\) by
\((E-E_{R;0,0})_+\) inside the second line while using that line for every
\(E>0\) would identify the two distinct cases in
\eqref{eq:s6cusp-native-count-threshold-values}: the zero value of the
positive part would retain the index pair \((0,0)\) even below the spectrum.

\begin{theorem}[Fixed-energy accumulation with the threshold retained]
\label{thm:s6cusp-native-fixed-energy-count}
For every fixed \(E>0\), as \(R\to\infty\),
\begin{align}
 N_{\eta,R}(E)
 &=\frac{\delta_R}{32\pi^2}
   \bigl(E-E_{R;0,0}\bigr)^2+O_E(R)
 \label{eq:s6cusp-native-fixed-energy-count-ground}\\
 &=\frac{E^2\delta_R}{32\pi^2}+O_E(R).
 \label{eq:s6cusp-native-fixed-energy-count-leading}
\end{align}
\end{theorem}

\begin{proof}
Equation \eqref{eq:s6cusp-native-ground-energy} and
\eqref{eq:s6cusp-native-frequency-asymptotics} give
\(E_{R;0,0}=14\pi R^{-1}+O(R^{-2})\).  Hence every fixed \(E>0\)
eventually satisfies \(E>E_{R;0,0}\), and the second branch of
\eqref{eq:s6cusp-native-count-piecewise} applies.  It counts the integer
points in the closed right triangle
\begin{equation}\label{eq:s6cusp-native-count-triangle}
 \left\{(x,y)\in\mathbb R_{\ge0}^2:
 2\mathcal B_{R,+}x+2\mathcal B_{R,-}y
 \le E-E_{R;0,0}\right\}.
\end{equation}
Its area, with both unequal intercepts retained, is
\begin{equation}\label{eq:s6cusp-native-count-triangle-area}
 \frac{(E-E_{R;0,0})^2}
 {8\mathcal B_{R,+}\mathcal B_{R,-}}.
\end{equation}
Attach the half-open unit square
\([n_+,n_++1)\times[n_-,n_-+1)\) to each counted lattice point.
The symmetric difference between their union and the triangle is contained
in a fixed-width neighborhood of the three boundary segments.  Its area is
bounded by a universal constant times the perimeter plus one.  The two axis
intercepts are
\begin{equation*}
 \frac{E-E_{R;0,0}}{2\mathcal B_{R,+}},
 \qquad
 \frac{E-E_{R;0,0}}{2\mathcal B_{R,-}},
\end{equation*}
and are \(O_E(R)\), so the lattice-point discrepancy is \(O_E(R)\).
Multiplication by six and the exact identity
\(\mathcal B_{R,+}\mathcal B_{R,-}=24\pi^2/\delta_R\)
prove \eqref{eq:s6cusp-native-fixed-energy-count-ground}.  Finally,
\begin{align*}
 \delta_R\left[(E-E_{R;0,0})^2-E^2\right]
 &=-2E\delta_RE_{R;0,0}
   +\delta_RE_{R;0,0}^2\\
 &=O_E(R),
\end{align*}
because \(\delta_R=R^2+O(R)\) and
\(E_{R;0,0}=14\pi R^{-1}+O(R^{-2})\).  This proves
\eqref{eq:s6cusp-native-fixed-energy-count-leading}.
\end{proof}

\subsection{The native topological split \(SU(2)\) saddle}
\label{subsec:s6cusp-native-split-su2}

Form the rank-two bundle
\begin{equation}\label{eq:s6cusp-native-split-bundle}
 \mathcal E_{\eta,R}:=
 \mathcal L_{\eta,R}\oplus\mathcal L_{\eta,R}^{-1}.
\end{equation}
The determinant trivialization is the evaluation pairing between the two
summands.  Retaining both Chern roots gives
\begin{align}
 c(\mathcal E_{\eta,R})
 &=(1+\eta_R)(1-\eta_R)=1-\eta_R^2,
 \label{eq:s6cusp-native-split-total-chern}\\
 c_1(\mathcal E_{\eta,R})&=0,
 &c_2(\mathcal E_{\eta,R})&=-\eta_R^2,
 &\int_{F_R}c_2(\mathcal E_{\eta,R})&=-12.
 \label{eq:s6cusp-native-split-c2}
\end{align}
Thus this split \(SU(2)\) bundle is not the trivial rank-two bundle: a bundle
isomorphism would pull back the zero second Chern class of the latter to the
nonzero class in \eqref{eq:s6cusp-native-split-c2}.

Put
\begin{equation}\label{eq:s6cusp-native-Hmag}
 H_{\mathrm{mag}}:=\operatorname{diag}(i,-i),
 \qquad
 \langle X,Y\rangle_{\mathfrak{su}(2)}
 :=-\frac12\operatorname{Tr}(XY),
\end{equation}
so that \(\lVert H_{\mathrm{mag}}\rVert=1\).  The direct-sum connection
\begin{equation}\label{eq:s6cusp-native-split-connection}
 A_{\eta,R}:=
 \nabla_{\eta,R}\oplus\nabla_{\eta,R}^{-1}
\end{equation}
has the exact curvature
\begin{equation}\label{eq:s6cusp-native-split-curvature}
 F_{A_{\eta,R}}=-2\pi\eta_RH_{\mathrm{mag}}.
\end{equation}
All coefficients of \(F_{A_{\eta,R}}\) in the covering coordinates are
constant, and both connection summands commute with
\(H_{\mathrm{mag}}\).  Therefore
\begin{equation}\label{eq:s6cusp-native-split-YM-equations}
 \dd_{A_{\eta,R}}F_{A_{\eta,R}}=0,
 \qquad
 \dd_{A_{\eta,R}}^*F_{A_{\eta,R}}=0.
\end{equation}
The first identity is also the Bianchi identity; the second follows directly
from the constant canonical coefficients and the flat covering metric.

Use throughout this subsection the action
\begin{equation}\label{eq:s6cusp-native-YM-action}
 \operatorname{YM}_R(A):=
 \frac1{2g_{\mathrm{YM}}^2}
 \int_{F_R}\lVert F_A\rVert^2\operatorname{vol}_{g_R}.
\end{equation}

\begin{theorem}[Exact action and exact topological excess]
\label{thm:s6cusp-native-split-action}
The Yang--Mills critical connection \(A_{\eta,R}\) satisfies
\begin{align}
 \operatorname{YM}_R(A_{\eta,R})
 &=\frac{\delta_R}{2g_{\mathrm{YM}}^2}
   \left(\mathcal B_{R,+}^2+\mathcal B_{R,-}^2\right)
 \label{eq:s6cusp-native-action-frequencies}\\
 &=\frac{2\pi^2}{g_{\mathrm{YM}}^2\delta_R}
   \left(p_R^2+36T_R^2+72m_R^2\right)
 \label{eq:s6cusp-native-action-coordinates}\\
 &=\frac{24\pi^2}{g_{\mathrm{YM}}^2}
   +\frac{2\pi^2(p_R+6T_R)^2}
   {g_{\mathrm{YM}}^2\delta_R}.
 \label{eq:s6cusp-native-action-excess}
\end{align}
Along the source radial cusp,
\begin{equation}\label{eq:s6cusp-native-action-cusp}
 \operatorname{YM}_R(A_{\eta,R})
 =\frac{74\pi^2}{g_{\mathrm{YM}}^2}
  +\frac{70\pi^2d_0}{g_{\mathrm{YM}}^2R}
  +O(R^{-2}).
\end{equation}
\end{theorem}

\begin{proof}
In the positive canonical planes, the two signed coefficients of \(\eta_R\)
have absolute values
\(\mathcal B_{R,+}/(2\pi)\) and
\(\mathcal B_{R,-}/(2\pi)\).  Equations
\eqref{eq:s6cusp-native-Hmag} and
\eqref{eq:s6cusp-native-split-curvature} therefore give
\begin{equation*}
 \lVert F_{A_{\eta,R}}\rVert^2
 =\mathcal B_{R,+}^2+\mathcal B_{R,-}^2.
\end{equation*}
The marked Euclidean volume is
\(|\det P_R|=\delta_R\), proving
\eqref{eq:s6cusp-native-action-frequencies}.  Since the signed canonical
coefficients are the eigenvalues of \(K_R/\delta_R\),
\begin{equation*}
 \kappa_{R,+}^2+\kappa_{R,-}^2
 =\operatorname{Tr}(K_R^2)
 =p_R^2+36T_R^2+72m_R^2.
\end{equation*}
Substitution proves \eqref{eq:s6cusp-native-action-coordinates}.  Direct
expansion, with \(p_R,T_R,m_R\) retained, gives
\begin{equation}\label{eq:s6cusp-native-action-coordinate-decomposition}
 p_R^2+36T_R^2+72m_R^2
 =12(6m_R^2-p_RT_R)+(p_R+6T_R)^2
 =12\delta_R+(p_R+6T_R)^2.
\end{equation}
This proves \eqref{eq:s6cusp-native-action-excess}.

For the radial expansion, equations
\eqref{eq:s6cusp-imaginary-identities} and
\eqref{eq:s6cusp-delta-two-term-expansion} retain
\begin{align*}
 p_R+6T_R
 &=5R+(d_0+5h_0)+O(e^{-2\pi R}),\\
 \delta_R
 &=R^2+(2h_0-d_0)R
   +(h_0^2-d_0h_0+6m_0^2)+O(Re^{-2\pi R}).
\end{align*}
Long division without deleting either displayed polynomial yields
\begin{equation*}
 \frac{(p_R+6T_R)^2}{\delta_R}
 =25+\frac{35d_0}{R}+O(R^{-2}).
\end{equation*}
Inserting this in \eqref{eq:s6cusp-native-action-excess} proves
\eqref{eq:s6cusp-native-action-cusp}.
\end{proof}

\subsection{Complete Coulomb-reduced Jacobi index of the native saddle}
\label{subsec:s6cusp-native-jacobi-index}

Let
\begin{equation}\label{eq:s6cusp-native-adjoint-root-splitting}
 \operatorname{ad}(\mathcal E_{\eta,R})_{\mathbb C}
 =\underline{\mathbb C}H_{\mathrm{mag}}
 \oplus\mathcal L_{\eta,R}^{2}E_+
 \oplus\mathcal L_{\eta,R}^{-2}E_-,
 \qquad
 [H_{\mathrm{mag}},E_\pm]=\pm2iE_\pm.
\end{equation}
The root-line Chern matrix is \(2E_\eta\), and therefore
\begin{equation}\label{eq:s6cusp-native-root-pfaffian}
 \left|\operatorname{Pf}(2E_\eta)\right|
 =4\left|\operatorname{Pf}(E_\eta)\right|=24.
\end{equation}
Applying the proved oscillator construction of
\cref{thm:s6cusp-native-four-dimensional-landau} to
\(\mathcal L_{\eta,R}^{2}\) gives scalar root levels
\begin{equation}\label{eq:s6cusp-native-root-scalar-levels}
 \lambda^{(0)}_{R;n_+,n_-}
 =2(2n_++1)\mathcal B_{R,+}
  +2(2n_-+1)\mathcal B_{R,-},
\end{equation}
each with complex multiplicity twenty-four.

On the complexified one-form Hilbert space define the full gauge-fixed
Jacobi operator with domain \(H^2\) by
\begin{equation}\label{eq:s6cusp-native-full-jacobi}
 (\mathcal J^{\mathrm{gf}}_{\eta,R}a)_\nu
 :=-D_{A_{\eta,R}}^\mu D_{A_{\eta,R},\mu}a_\nu
   -2[F_{A_{\eta,R},\nu\mu},a_\mu].
\end{equation}
This is a symmetric elliptic Laplace-type operator on a smooth bundle over
the compact torus, so its \(H^2\) realization is self-adjoint and has compact
resolvent.  Define the closed Coulomb Hilbert subspace by
\begin{equation}\label{eq:s6cusp-native-Coulomb-Hilbert-space}
 \mathscr C_{\eta,R}:=
 \left(\overline{\operatorname{im}
  (\dd_{A_{\eta,R}}:H^1(\operatorname{ad}\mathcal E_{\eta,R})
   \longrightarrow L^2(T^*F_R\otimes\operatorname{ad}\mathcal E_{\eta,R}))}
 \right)^\perp.
\end{equation}
For an \(H^1\) one-form, membership is equivalent to the weak equation
\(\dd_{A_{\eta,R}}^*a=0\).

The Coulomb space reduces \(\mathcal J^{\mathrm{gf}}_{\eta,R}\).  To prove
this, let \(\mathcal L^{\mathrm{YM}}_{\eta,R}\) denote the ungauged
Yang--Mills Hessian differential at the critical connection.  Gauge
covariance of
\(A\mapsto\dd_A^*F_A\) gives, after differentiating the path
\(\exp(t\phi)\cdot A_{\eta,R}\) at \(t=0\),
\begin{equation}\label{eq:s6cusp-native-gauge-orbit-hessian-zero}
 \mathcal L^{\mathrm{YM}}_{\eta,R}\dd_{A_{\eta,R}}\phi=0.
\end{equation}
The gauge-fixed operator is
\(\mathcal J^{\mathrm{gf}}_{\eta,R}
=\mathcal L^{\mathrm{YM}}_{\eta,R}
 +\dd_{A_{\eta,R}}\dd_{A_{\eta,R}}^*\).
Hence, on smooth zero-forms,
\begin{equation}\label{eq:s6cusp-native-jacobi-gradient-intertwiner}
 \mathcal J^{\mathrm{gf}}_{\eta,R}\dd_{A_{\eta,R}}
 =\dd_{A_{\eta,R}}\Delta_{A_{\eta,R},0},
 \qquad
 \Delta_{A_{\eta,R},0}:=
 \dd_{A_{\eta,R}}^*\dd_{A_{\eta,R}}.
\end{equation}
For smooth \(a\in\mathscr C_{\eta,R}\) and smooth \(\phi\), self-adjointness
and \eqref{eq:s6cusp-native-jacobi-gradient-intertwiner} give
\begin{align*}
 \left\langle \dd_{A_{\eta,R}}\phi,
 \mathcal J^{\mathrm{gf}}_{\eta,R}a\right\rangle
 &=\left\langle
 \mathcal J^{\mathrm{gf}}_{\eta,R}\dd_{A_{\eta,R}}\phi,a
 \right\rangle\\
 &=\left\langle
 \dd_{A_{\eta,R}}\Delta_{A_{\eta,R},0}\phi,a
 \right\rangle=0.
\end{align*}
Density and closedness prove invariance of \(\mathscr C_{\eta,R}\); the
orthogonal complement is invariant as well.  The restricted operator
\begin{equation}\label{eq:s6cusp-native-Coulomb-Jacobi}
 \mathcal J^{\mathrm C}_{\eta,R}:=
 \mathcal J^{\mathrm{gf}}_{\eta,R}\big|_{
 \operatorname{Dom}(\mathcal J^{\mathrm{gf}}_{\eta,R})
 \cap\mathscr C_{\eta,R}}
\end{equation}
is therefore self-adjoint with discrete spectrum.

\begin{theorem}[Three negative bands and exact real Morse index]
\label{thm:s6cusp-native-jacobi-index}
There is \(R_4\) such that every \(R\ge R_4\) has
\begin{equation}\label{eq:s6cusp-native-ratio-five-seven}
 5<\frac{\mathcal B_{R,+}}{\mathcal B_{R,-}}<7.
\end{equation}
For those \(R\), the complete negative spectrum of
\(\mathcal J^{\mathrm C}_{\eta,R}\) consists of the three distinct
eigenvalues
\begin{align}
 \nu_{R,n}
 &:=2\left[(2n+1)\mathcal B_{R,-}-\mathcal B_{R,+}\right]
 \label{eq:s6cusp-native-negative-bands}\\
 &=\frac{4\pi}{\delta_R}
 \left[n\Xi_R-(n+1)S_R\right],
 \qquad n=0,1,2.
 \label{eq:s6cusp-native-negative-bands-coordinates}
\end{align}
Every one of these eigenspaces has real dimension forty-eight.  Consequently
\begin{equation}\label{eq:s6cusp-native-Morse-index}
 \operatorname{index}_{\mathbb R}
 (\mathcal J^{\mathrm C}_{\eta,R})=3\cdot48=144.
\end{equation}
Their leading cusp terms, together with the next positive member of the same
branch, are
\begin{align}
 \nu_{R,0}&=-\frac{20\pi}{R}+O(R^{-2}),&
 \nu_{R,1}&=-\frac{12\pi}{R}+O(R^{-2}),\notag\\
 \nu_{R,2}&=-\frac{4\pi}{R}+O(R^{-2}),&
 \nu_{R,3}&= \frac{4\pi}{R}+O(R^{-2}).
 \label{eq:s6cusp-native-negative-band-cusp}
\end{align}
\end{theorem}

\begin{proof}
The frequency ratio tends to six by
\eqref{eq:s6cusp-native-frequency-asymptotics}; its strict containment in
\((5,7)\) follows for all \(R\) beyond one fixed \(R_4\).

Use the exact canonical coordinates supplied by
\eqref{eq:s6cusp-native-orthogonal-map}.  On the \(E_+\) root line, the
spin-curvature part of \eqref{eq:s6cusp-native-full-jacobi} acts on the
\((\dd X_+,\dd Y_+)\) coefficients by
\begin{equation}\label{eq:s6cusp-native-strong-spin-matrix}
 \mathsf S_{R,+}
 =4\mathcal B_{R,+}
 \begin{pmatrix}0&i\\-i&0\end{pmatrix},
\end{equation}
and on the \((\dd X_-,\dd Y_-)\) coefficients by the same matrix with
\(\mathcal B_{R,+}\) replaced by \(\mathcal B_{R,-}\).  Direct
multiplication gives
\begin{align*}
 \mathsf S_{R,+}\binom{1}{i}
 &=-4\mathcal B_{R,+}\binom{1}{i},&
 \mathsf S_{R,+}\binom{1}{-i}
 &=4\mathcal B_{R,+}\binom{1}{-i}.
\end{align*}
Thus the four constant circular one-form polarizations shift the scalar
levels \eqref{eq:s6cusp-native-root-scalar-levels} by
\(-4\mathcal B_{R,+},+4\mathcal B_{R,+},
-4\mathcal B_{R,-},+4\mathcal B_{R,-}\), respectively.

For the negative strong-plane polarization the complete level is
\begin{equation}\label{eq:s6cusp-native-strong-negative-full-level}
 2\left[(2n_+-1)\mathcal B_{R,+}
       +(2n_-+1)\mathcal B_{R,-}\right].
\end{equation}
When \(n_+\ge1\), both coefficients in brackets are positive.  When
\(n_+=0\), negativity is equivalent to
\begin{equation*}
 2n_-+1<\frac{\mathcal B_{R,+}}{\mathcal B_{R,-}}.
\end{equation*}
Under \eqref{eq:s6cusp-native-ratio-five-seven}, the odd integers satisfying
this inequality are exactly \(1,3,5\), so
\(n_-=0,1,2\).  The next odd integer is seven, and the strict upper bound in
\eqref{eq:s6cusp-native-ratio-five-seven} makes the \(n_-=3\) level
positive.  The weak-plane negative polarization has minimum
\begin{equation*}
 2(\mathcal B_{R,+}-\mathcal B_{R,-})>0,
\end{equation*}
the two positive spin shifts are positive, and the Cartan sector is the
nonnegative ordinary Hodge Laplacian.  This exhausts the Cartan summand, both
root summands, and all four one-form polarizations.

It remains to show that restricting to Coulomb gauge removes none of the
negative spectrum.  On \(\mathcal L_{\eta,R}^{2}\), put
\begin{equation*}
 \mathfrak b_+:=
 \frac{\Pi_{X_+}+i\Pi_{Y_+}}
 {\sqrt{4\mathcal B_{R,+}}}.
\end{equation*}
Every \(n_+=0\) scalar state \(\phi_{0,n_-}\) is killed by
\(\mathfrak b_+\).  Its negative-polarization one-form obeys the exact
divergence identity
\begin{align}
 \dd_{A_{\eta,R}}^*
 \left[\phi_{0,n_-}(\dd X_++i\dd Y_+)E_+\right]
 &=-i\sqrt{4\mathcal B_{R,+}}
   (\mathfrak b_+\phi_{0,n_-})E_+\notag\\
 &=0.
 \label{eq:s6cusp-native-negative-Coulomb-equation}
\end{align}
The Hermitian-adjoint \(E_-\) modes satisfy the conjugate equation.  Hence
the complete negative spectral subspace of the full operator lies in the
reducing Coulomb subspace.

Equation \eqref{eq:s6cusp-native-root-pfaffian} and the root oscillator
proof give complex dimension twenty-four for each \(E_+\) scalar joint
level.  Hermitian adjunction identifies the \(E_-\) level with its complex
conjugate.  If \(v_1,\ldots,v_{24}\) is a complex basis in the \(E_+\)
level, then
\begin{equation*}
 v_j-v_j^\dagger,
 \qquad i(v_j+v_j^\dagger),
 \qquad 1\le j\le24,
\end{equation*}
are forty-eight real linearly independent anti-Hermitian eigenforms and
span the real root eigenspace.  This proves the multiplicity and
\eqref{eq:s6cusp-native-Morse-index}.  Substitution of the exact frequencies
gives \eqref{eq:s6cusp-native-negative-bands-coordinates}; substitution of
their retained cusp expansions gives
\eqref{eq:s6cusp-native-negative-band-cusp}.
\end{proof}

The action excess and the lowest Jacobi branch satisfy the exact identity
\begin{align}
 \nu_{R,0}
 &=-2(\mathcal B_{R,+}-\mathcal B_{R,-})
  =-\frac{4\pi(p_R+6T_R)}{\delta_R},
 \label{eq:s6cusp-native-lowest-instability}\\
 \operatorname{YM}_R(A_{\eta,R})
 -\frac{24\pi^2}{g_{\mathrm{YM}}^2}
 &=\frac{\delta_R}{8g_{\mathrm{YM}}^2}\nu_{R,0}^2.
 \label{eq:s6cusp-native-action-instability}
\end{align}
Indeed, the first equality follows from the exact frequency difference, and
substitution into \eqref{eq:s6cusp-native-action-excess} proves the second.
The spin-one shift and its sign agree with the homogeneous-background
operator \(K_\pm=-D^2\pm2gB\) and its lowest negative branch in
\cite[eqs.~(1)--(3)]{Oxman2010}; the compact bundle, the twenty-four theta
states, the Coulomb equation, and the index \(144\) are proved above rather
than imported from that comparison.

\subsection{One decomposable native operator with continuous threshold zero}
\label{subsec:s6cusp-native-direct-integral}

The marking supplies an exact common Hilbert-space model.  Let
\(\mathbb T_y^4:=\mathbb R_y^4/\mathbb Z^4\), let \(\mathcal L_\eta\) be
the fixed line bundle with clutch factors
\eqref{eq:s6cusp-native-line-clutch-one}--
\eqref{eq:s6cusp-native-line-clutch-two}, and define
\begin{equation}\label{eq:s6cusp-native-marking-diffeomorphism}
 \varphi_R:\mathbb T_y^4\longrightarrow F_R,
 \qquad
 \varphi_R([y]):=[P_Ry].
\end{equation}
Its inverse is \([x]\mapsto[P_R^{-1}x]\), so it is a diffeomorphism and
identifies \(\varphi_R^*\mathcal L_{\eta,R}\) with \(\mathcal L_\eta\).
The pulled-back metric and volume form are
\begin{equation}\label{eq:s6cusp-native-pulled-metric-volume}
 \varphi_R^*g_R=(\dd y)^{\mathsf T}P_R^{\mathsf T}P_R\dd y,
 \qquad
 \varphi_R^*\operatorname{vol}_{g_R}
 =\delta_R\dd y^1\dd y^2\dd y^3\dd y^4.
\end{equation}
Consequently
\begin{equation}\label{eq:s6cusp-native-marking-unitary}
 U_R:L^2(F_R,\mathcal L_{\eta,R};\operatorname{vol}_{g_R})
 \longrightarrow
 \mathscr K:=L^2(\mathbb T_y^4,\mathcal L_\eta;\dd^4y),
 \qquad
 (U_Rs)(y):=\delta_R^{1/2}s(P_Ry)
\end{equation}
is unitary.  Put
\begin{equation}\label{eq:s6cusp-native-fixed-fibre-operator}
 \widehat H_R:=U_RH_{\eta,R}U_R^{-1},
 \qquad
 \operatorname{Dom}(\widehat H_R)
 =H^2(\mathbb T_y^4,\mathcal L_\eta).
\end{equation}
The metric matrix \(P_R^{\mathsf T}P_R\), its inverse, and the fixed
connection coefficients make \(R\mapsto\widehat H_R\) a smooth family from
the common \(H^2\) domain to \(\mathscr K\).

Fix \(R_0\) in the regular radial range and define
\begin{align}
 \mathscr H_\eta
 &:=L^2([R_0,\infty),\dd R;\mathscr K),
 \label{eq:s6cusp-native-direct-integral-Hilbert}\\
 (\mathscr M_\eta\Psi)(R)&:=\widehat H_R\Psi(R),
 \label{eq:s6cusp-native-direct-integral-action}\\
 \operatorname{Dom}(\mathscr M_\eta)
 &:=\left\{\Psi\in\mathscr H_\eta:
 \begin{array}{l}
 \Psi(R)\in H^2(\mathbb T_y^4,\mathcal L_\eta)
 \text{ for almost every }R,\\
 R\mapsto\widehat H_R\Psi(R)\text{ is measurable, and}\\
 \displaystyle\int_{R_0}^{\infty}
 \lVert\widehat H_R\Psi(R)\rVert_{\mathscr K}^2\dd R<\infty
 \end{array}\right\}.
 \label{eq:s6cusp-native-direct-integral-domain}
\end{align}

\begin{theorem}[Self-adjointness and continuous essential threshold]
\label{thm:s6cusp-native-direct-integral-threshold}
The operator \(\mathscr M_\eta\) is self-adjoint and nonnegative, and
\begin{equation}\label{eq:s6cusp-native-direct-integral-spectrum-zero}
 0\in\operatorname{Spec}_{\mathrm{ess}}(\mathscr M_\eta)
 \cap\operatorname{Spec}_{\mathrm c}(\mathscr M_\eta),
 \qquad
 \ker\mathscr M_\eta=\{0\}.
\end{equation}
\end{theorem}

\begin{proof}
The common-domain smoothness in
\eqref{eq:s6cusp-native-fixed-fibre-operator}, the local elliptic estimate,
and the resolvent identity imply local norm continuity of
\(R\mapsto(\widehat H_R\pm i)^{-1}\).  In particular the resolvent field is
strongly measurable and has norm at most one.  Define the bounded operators
on \(\mathscr H_\eta\)
\begin{equation*}
 (\mathscr R_\pm\Psi)(R):=
 (\widehat H_R\pm i)^{-1}\Psi(R).
\end{equation*}
For \(\Phi=\mathscr R_\pm\Psi\),
\begin{equation*}
 \widehat H_R\Phi(R)=\Psi(R)\mp i\Phi(R),
\end{equation*}
so \(\Phi\) belongs to
\(\operatorname{Dom}(\mathscr M_\eta)\) and
\((\mathscr M_\eta\pm i)\Phi=\Psi\).  Fibrewise integration proves
symmetry of \(\mathscr M_\eta\); surjectivity of both
\(\mathscr M_\eta\pm i\) proves self-adjointness.  Equation
\eqref{eq:s6cusp-native-ground-lower-bound} proves nonnegativity.

The rank-six ground projection is isolated at each finite \(R\) by the
strict gap \(2\min\{\mathcal B_{R,+},\mathcal B_{R,-}\}\).  A local Riesz
contour and the norm-continuous resolvent therefore give a smooth rank-six
projection field.  These ranges form a smooth vector bundle over the
interval \([R_0,\infty)\); the interval is contractible, so choose a smooth
unit section \(R\mapsto\phi_R\) with
\begin{equation}\label{eq:s6cusp-native-ground-section}
 \widehat H_R\phi_R=E_{R;0,0}\phi_R,
 \qquad \lVert\phi_R\rVert_{\mathscr K}=1.
\end{equation}
For every integer \(j>R_0\), put
\begin{equation}\label{eq:s6cusp-native-Weyl-sequence}
 \Psi_j(R):=\mathbf1_{[j,j+1]}(R)\phi_R.
\end{equation}
The vectors are orthonormal; hence they converge weakly to zero.  Equations
\eqref{eq:s6cusp-native-ground-section} and
\(E_{R;0,0}=14\pi R^{-1}+O(R^{-2})\) give
\begin{equation}\label{eq:s6cusp-native-Weyl-residual}
 \lVert\mathscr M_\eta\Psi_j\rVert^2
 =\int_j^{j+1}E_{R;0,0}^2\dd R
 =O(j^{-2})\longrightarrow0.
\end{equation}
Thus \((\Psi_j)\) is a singular Weyl sequence at zero, proving
\(0\in\operatorname{Spec}_{\mathrm{ess}}(\mathscr M_\eta)\).

If \(\mathscr M_\eta\Psi=0\), then
\(\widehat H_R\Psi(R)=0\) for almost every \(R\).  The strict fibrewise
bound
\(\widehat H_R\ge E_{R;0,0}I>0\) forces \(\Psi(R)=0\) almost everywhere,
so the kernel is zero.  A self-adjoint operator has no residual spectrum;
zero is therefore in the continuous spectrum, which proves
\eqref{eq:s6cusp-native-direct-integral-spectrum-zero}.
\end{proof}

The exact morphism from the radial family to this operator is the marking
unitary \eqref{eq:s6cusp-native-marking-unitary} followed by the decomposable
action \eqref{eq:s6cusp-native-direct-integral-action}.  There is no
derivative in the \(R\)-direction: a horizontal kinetic completion would
contain an additional differential expression such as \(-\partial_R^2\)
together with a specified measure and boundary domain.  Hence
\cref{thm:s6cusp-native-direct-integral-threshold} proves a continuous
threshold for the displayed decomposable magnetic operator, not the
spectrum of an unconstructed horizontal completion or of a reconstructed
interacting Yang--Mills Hamiltonian.
